\documentclass{article}
\usepackage{pathfinder-readable}

\title{A cooperative route through asymmetric contracting boards}

\begin{document}
\maketitle

\begin{abstract}
This note connects a theory of how to make agents report and act as intended with a game in which
language-model agents negotiate chip transfers. We asked whether poor cooperation on boards that
favour one player reflects a lack of a good outcome, a limit of the contracts, or a failure to find
and carry out suitable terms. A checked route and a five-tile contract let both players beat their
separate outside options on every asymmetric board generated by the game's stated rule. A score
identity shows how much coverage the stronger player must actually receive for this to happen. The
thread ended at DRAFT because the witness establishes an attainable outcome, while the observed
failures still lack a tested explanation and no full bargaining equilibrium has been shown.
\end{abstract}

\section{Introduction}

Bergemann, Koh and Morris's \emph{Mechanism Design for Alignment and Control}, called Q here,
studies a designer who asks agents to act on their behalf while their preferences and capabilities
may be hidden. It distinguishes an action that is physically available from a rule that gives an
agent reason to report accurately and obey a recommendation. Its revelation results concern reports
followed by a one-shot action. A reward can also respond only to distinctions the evaluator can
observe \cite{Q}.

Wyse and colleagues' \emph{Commitment To Cooperation With Self-Negotiated Contracts}, called P,
tests language-model agents in a two-player board game. Red and Blue each try to reach a goal by
spending chips that match the colours of the squares they enter. They may trade chips, promise to
pay for the other's moves, or negotiate a contract. On an asymmetric board Red can finish without
help, whereas Blue needs a Red chip. P measures whether both players score strictly more than they
could without interacting \cite{P}.

The connection was a diagnostic question. When both players rarely improve on an asymmetric board,
is a mutually beneficial route absent, are the allowed agreements too weak, or do the agents fail to
choose and use one? Q's distinction between feasibility and incentives helps frame the question. Its
one-shot results do not settle the bargaining and movement decisions in P's sequential game.

\section{What was done}

The peers first examined P's voluntary Pay-for-Partner promises and its point contracts. One spent
chip costs a player who finishes five score points. They calculated when a promised payment on the
partner's completion could make covering a tile worthwhile. In parallel, they examined P's
difference between contracts written as tile-by-tile instructions and contracts interpreted from
natural language. That comparison raised a question about which tiles were agreed and which
attempted moves a judge approved, but it did not identify the cause of the performance difference.

The search then moved to the boards themselves. Ada enumerated all \(3{,}432\) ways to colour the
fourteen non-green squares with seven Red and seven Blue squares under P's stated generation rule.
Emmy reran the oracle and checked whether longer simple routes could reduce Blue's need for Red
chips. They next found an immediately executed trade that gives both players a good route. Finally,
Ada wrote a reciprocal programmatic tile contract for the Pay-for-Partner setting and derived a
score test based on coverage that actually occurs. Emmy checked the trade against P's interface; the
contract and score calculations are recorded in \url{../ada/asymmetric_oracle.py} and
\url{../ada/bottleneck_contracts.md}.

\section{Results}

Among those balanced boards, \(396\) have an independent Red player and a dependent Blue player.
Each has a six-move Red-only route. On a shortest Blue route, let \(k\) be the number of Red squares
among the five coloured moves; the sixth move enters the green goal. The check found \(k=1\) on
\(314\) boards and \(k=2\) on \(82\). Emmy's enumeration of all simple routes found that taking a
longer route never reduces this minimum Red-chip need; see \url{../emmy/observability-frontier.md}.

For each of these boards, both players can take the same shortest Blue route. A programmatic tile
contract makes Red pay for Blue's move onto each of the \(k\) Red squares and Blue pay for Red's
move onto each of the other \(5-k\) Blue squares. These are five tile-specific clauses of the form P
permits. Each player pays for their own moves on squares of their own colour and for their own entry
into the green goal. Transfers under an accepted tile contract occur automatically when the covered
move is attempted, provided the giver has the chip \cite{P}.

Each player starts with sixteen chips. P awards a finisher twenty points plus five points for each
chip left. If \(r_R\) and \(r_B\) are Red's and Blue's final scores, the route gives
\[
(r_R,r_B)=(95-10k,\,45+10k).
\]
For \(k=1\), this is \((85,55)\); for \(k=2\), it is \((75,65)\). Both exceed P's no-interaction
scores \((70,0)\), and the total is the maximum joint score, \(140\). The oracle checked the
inventories and scores for every eligible board. This is a witness for a feasible contract and path
if the agreement is accepted and the route is followed. It does not show that agents will propose or
accept it.

The score rule gives a second, more general test. Suppose Red finishes after \(L_R\) moves, Blue has
paid for \(c\) of Red's moves, and Red has paid for \(d\) of Blue's. Assume there are no other chip
trades or point transfers. Since a shortest route has six moves, Red's score is
\[
r_R=70+5\bigl[c-d-(L_R-6)\bigr].
\]
Red therefore strictly beats its outside score precisely when \(c-d>L_R-6\). On a shortest route,
Red must receive more coverage than it gives. For the shared route above, \(c=5-k\), \(d=k\), and
\(L_R=6\), so the condition holds for either value of \(k\). The quantities are delivered chip
payments, not merely promised clauses. P already measures net tiles promised; comparing that measure
with delivered coverage and route length would test where this condition fails \cite{P}.

P reports a mean rate of \(0.06\pm0.03\) for both players beating baseline with Pay-for-Partner and
programmatic tile contracts on asymmetric boards, alongside mean contract acceptance of
\(0.80\pm0.03\) \cite{P}. The all-board witness gives a feasible rate of one under its prescribed
terms and routes. Those aggregate observations cannot say whether the shortfall arose in proposing
terms, accepting them, choosing moves, or making later payments. The peers also found a
regular-trading witness: trade \(k\) Red chips for \(k+1\) Blue chips before movement, then take
six-move routes. This gives \((75,65)\) on every eligible board, while P reports a mean
both-beat-baseline rate of \(0.20\) for regular trading without contracts \cite{P}. That comparison
likewise establishes an attainable outcome, not the result of voluntary bargaining.

\section{Objections and limits}

An early explanation was that P's point contracts cannot promise more than twenty points. At a fixed
coverage decision, let \(x\) be a payment to the covering player if the partner finishes, and let
\(p_1\) and \(p_0\) be that player's perceived probabilities of partner completion after coverage
and refusal under the same continuation assumptions. If the covering player finishes in either case
and nothing else changes, coverage is weakly worthwhile exactly when \(x(p_1-p_0)\geq5\). If the
probability difference is below one quarter, the twenty-point cap cannot supply this incentive.
Emmy derived a related limit for a bonus paid on an observable signal in
\url{../emmy/observability-frontier.md}. These are local incentive calculations, not equilibrium results.
The cap alone does not explain the present board witness: suitable short routes need only one or two
Red-chip coverages. A proposed point-contract construction also depended on an unconfirmed one-sided
Pay-for-Partner offer, which the reciprocal tile contract avoids.

The natural-language contract comparison raised another objection. P reports different completion
rates for programmatic and natural-language tile contracts on mutually dependent boards, even though
their average numbers of covered tiles are similar \cite{P}. Similar counts do not show that the
same squares were covered. The judge audit counted errors among approved moves and so could not
measure wrongly rejected eligible moves. Emmy showed that two formats with the same accepted terms,
information, and transfer decisions at every history would induce the same game; see
\url{../emmy/contract-representation-note.md}. Whether P's formats meet those conditions remains
untested. Neither this comparison nor the asymmetric-board aggregates isolate a cause of failure.

\section{Why the thread stopped}

The verifier left the thread at DRAFT. The ledger and peer checks support the all-board
contract-and-path witness and the score identity, but they do not establish that agents would
negotiate the contract, obey it at every later decision, or fail for any particular measured reason.
The smallest next test is to supply the verified reciprocal contract as already accepted on the
asymmetric boards, let agents move, and record route length and delivered coverage. Success would
show that agents can execute these supplied terms, leaving their discovery and acceptance as
separate questions. Failure would show an execution obstacle even with those terms supplied.

\bibliographystyle{plainurl}
\bibliography{references}
\end{document}
